При пороге 0,1 Laya с описаниями Qwen3-1.7B дала AMI 0,374, с описаниями Qwen3.5-4B --- 0,367. Центры, пять представителей каждого кластера и тестовые письма совпадают. После отдельного подбора порога для Laya ($\tau=1$) и GPT-6-Luna ($\tau=0,3$) тестовая AMI составила 0,548 и 0,547 соответственно.

В таблице для BigARTM приведена заранее заданная инициализация 11. Для проверки устойчивости повторили обучение с инициализациями 23 и 47. У PLSA диапазон тестовой AMI по трём запускам составил 0,042--0,103, у сглаженной модели --- 0,072--0,110. Во всех шести случаях на тестовых документах представлены все 20 тем.

При штатном пороге GPT-6-Luna вернула Other для 140 тестовых писем, подтвердила исходную тему для 14 и сменила её для 6. Большая доля отказов объясняет близость результата к исходному разбиению.
